\documentclass[UTF8,12pt,a4paper]{ctexart}

\usepackage{amsmath,amssymb}
\usepackage{booktabs}
\usepackage{geometry}
\usepackage{graphicx}
\usepackage[hidelinks]{hyperref}
\geometry{left=2.5cm,right=2.5cm,top=2.5cm,bottom=2.5cm}
\newcommand{\keywords}[1]{\par\noindent\textbf{关键词：}#1}

% 国赛排版基线（无官方模板时）：一级标题黑体四号居中、中文数字编号；
% 二三级标题小四加粗左对齐；正文宋体小四、首行缩进 2 字符、约 1.5 倍行距。
\ctexset{
  section = {
    format = \zihao{4}\heiti\centering,
    name = {,、},
    number = \chinese{section},
  },
  subsection = {
    format = \zihao{-4}\heiti\raggedright,
  },
  subsubsection = {
    format = \zihao{-4}\bfseries\raggedright,
  },
}
\linespread{1.3}
\setlength{\parindent}{2em}

% 这是无官方 LaTeX 模板时的构建基线，不替代当届官方规则。
\title{LATEX_TEMPLATE_TITLE}
\author{}
\date{}

\begin{document}
\maketitle

\begin{abstract}
LATEX_TEMPLATE_ABSTRACT
\keywords{LATEX_TEMPLATE_KEYWORDS}
\end{abstract}

\section{问题重述}
LATEX_TEMPLATE_BODY

\section{问题分析}
LATEX_TEMPLATE_BODY

\section{模型假设}
LATEX_TEMPLATE_BODY

\section{符号说明}
LATEX_TEMPLATE_BODY

% 主体章按问题联动：每问一节（模型建立→模型求解→结果与分析），
% 禁止"先全部建模、再全部求解"两段式；公共底座前置不超过 1.5 页。
\section{模型的建立与求解}
LATEX_TEMPLATE_BODY

\section{模型评价}
LATEX_TEMPLATE_BODY

\bibliographystyle{plain}
\bibliography{references}

% 附录三大件：索引清单（支撑文件名+功能注释）/ 答案全精度大表 / 按问题分文件核心代码。
% \appendix
% \section{支撑材料清单}
% \section{问题求解答案}
% \section{程序代码}
\end{document}
